\section{Dificultades del diseño del producto: arranque en frío, doble arranque en frío y proporción de personas reales}

Elys no es un producto que funcione por sí solo con que el modelo sea lo bastante capaz. Debe resolver dos arranques en frío al mismo tiempo: el primero es el arranque en frío de usuarios, es decir, cómo lograr que suficientes personas entren en la red; el segundo es el arranque en frío de context, es decir, cómo lograr que los usuarios estén dispuestos a aportar información sobre sí mismos en cantidad suficiente, auténtica y utilizable. Solo si los usuarios y el context arrancan a la vez puede el doble digital generar conexiones reales.

\subsection{Doble arranque en frío}

\noindent\begin{longtable}{p{0.22\textwidth}p{0.34\textwidth}p{0.35\textwidth}}
\toprule
Tipo de arranque en frío & Problema central & Soluciones posibles \\
\midrule
Arranque en frío de usuarios & Sin suficientes personas reales, la interacción entre dobles parece girar en vacío & Arranque en círculos pequeños, acceso por invitación, comunidades semilla, difusión basada en escenarios de uso. \\
Arranque en frío de context & Los usuarios no quieren o no saben entregar context de alta calidad & Onboarding de baja fricción, beneficio inmediato, autorización revocable, guía con ejemplos. \\
Arranque en frío de confianza & A los usuarios les preocupa que la empresa controle su privacidad o que el doble los represente mal & Almacenamiento local, permisos por niveles, auditoría de acciones, confirmación de acciones sensibles. \\
Arranque en frío de la red & Hay personas pero no relaciones; hay contenido pero no conexiones & La IA rompe primero el hielo y después guía a las personas reales a confirmar e interactuar. \\
\bottomrule
\end{longtable}

\begin{importantbox}{La esencia del doble arranque en frío}
Los productos sociales tradicionales arrancan en frío sobre todo usuarios y contenido; la IA social además tiene que arrancar en frío el context. Sin context, la IA no se parece a ti; sin personas reales, la actividad de la IA no tiene sentido.
\end{importantbox}

\subsection{Por qué importa la proporción de personas reales}

La métrica estrella polar de Elys es la tasa de conexiones entre personas reales, o la proporción de acciones realizadas por personas reales. Este punto es decisivo: la IA puede reducir la barrera para interactuar, pero no puede hacer que en la red solo quede IA. Si cada vez que el usuario abre la aplicación solo ve respuestas de IA, se decepcionará, porque el sentido de lo social sigue proviniendo de las personas reales. El papel de la IA debe ser amplificar las conexiones entre personas reales, no sustituirlas.

\begin{warningbox}{La «trampa del bullicio» en la IA social}
Si el producto cuenta como crecimiento los comentarios, los «me gusta» y las publicaciones hechos por IA, los datos pueden verse bien a corto plazo, pero los usuarios descubrirán pronto que no hay personas reales. Lo que más teme una red social es la prosperidad ficticia, porque destruye la confianza.
\end{warningbox}

\subsection{Resumen de la sección}

La dificultad del producto Elys no está en «si puede generar conversaciones», sino en si puede poner en marcha al mismo tiempo a los usuarios, el context, la confianza y la red de personas reales. La tasa de conexiones entre personas reales es una métrica más cercana a lo esencial que el número de mensajes.
